% ============================================================================
%  front.tex —— 标题页、摘要、译者说明、目录、引言
% ============================================================================

\begin{titlepage}
  \centering
  \vspace*{2.2cm}
  {\Huge\bfseries 迈向同调镜像对称\par}
  \vspace{0.9cm}
  {\Large Towards Homological Mirror Symmetry\par}
  \vspace{2.2cm}
  {\large Alessandro Imparato\quad 著\par}
  \vspace{0.4cm}
  {\large 苏黎世联邦理工学院（ETH Zürich）数学系\par}
  \vspace{2.6cm}
  {\large 中译本\par}
  \vspace{0.3cm}
  {\normalsize 底本：arXiv:2108.03931v2 [math.SG]，139 页\par}
  \vfill
  {\normalsize 二〇二六年\par}
\end{titlepage}

% 前置部分用小写罗马数字页码，正文从阿拉伯数字 1 重新开始
\pagenumbering{roman}
\pagestyle{plain}

% ---------------------------------------------------------------------------
% 摘要
% ---------------------------------------------------------------------------
\section*{摘要}
\addcontentsline{toc}{section}{摘要}

本文是关于 Kontsevich 同调镜像对称猜想 A-侧的导论性文章。我们首先自足地研究
$A_\infty$-范畴及其同调代数，随后专门讨论 Fukaya 范畴，着重介绍其背后的 Floer 理论
基础以及其中的几何特征。最后，我们把这一理论置于镜像对称的语境中，逐步走向它的主要
预言。

本文是作者的硕士学位论文，由 William John Merry 教授指导，为获得苏黎世联邦理工学院
（ETH Zürich）数学硕士学位而提交。它应当对研究生而言是可读的。

\vspace{0.8em}
\noindent\textbf{关键词：}$A_\infty$-范畴，分裂封闭导出范畴，拉格朗日子流形，拉格朗日
Floer 上同调，Fukaya 范畴，镜像对称，同调镜像对称。

% ---------------------------------------------------------------------------
% 译者说明
% ---------------------------------------------------------------------------
\section*{译者说明}
\addcontentsline{toc}{section}{译者说明}

\begin{enumerate}[leftmargin=2.2em,itemsep=0.35em]
  \item \textbf{底本。}译文依据 arXiv:2108.03931v2（2023 年 6 月 21 日修订版，139 页，
    15 幅图）。原文页码在译文图题末尾以"（原书 p.N）"标出，正文交叉引用一律改用
    \LaTeX{} 标签，与原书编号保持一致（章内连续编号，如"定理 3.3"）。
  \item \textbf{术语。}采用中文学界通行译名（如 sheaf 译为层，derived category 译为导出
    范畴，triangulated 译为三角的，holonomy 译为和乐群）；尚无稳定译名者保留英文（如
    wrapped、stop、thimble）。关键术语首次出现时括注英文。书末附译者增补的
    《英中术语对照表》。
  \item \textbf{公式。}全部公式按原书页面逐符号转写，未作改写；原书疑似笔误一律照印，
    并在紧邻处以 LaTeX 注释说明，绝不悄悄修正。
  \item \textbf{批注。}源码中散布形如
    \verb|% [Tutor Note: 【易错】……]| 的注释，分【易错】（符号、指标、约定）、
    【含糊】（作者写"显然/可以证明"而省略的论证）、【勘误】（原书笔误）、
    【背景】（一两句直觉）四类。注释不进入排版结果，阅读 PDF 时不可见，需在源码中查看。
  \item \textbf{图表。}原书 15 幅位图由 PDF 无损提取后以 \verb|\includegraphics| 插入，
    图题译为中文并标注原书页码；正文中的交换图原为 TikZ 矢量图，译本用 \verb|tikz-cd|
    重排为可编辑源码，未转为位图。
  \item \textbf{节略。}原书书末索引（原书 pp.134–136）属于检索装置而非正文内容，译本未逐条
    复制，改以译者增补的术语对照表代之。
  \item \textbf{编译。}本文件须用 \XeLaTeX{} 编译（中文由 \texttt{ctex} 处理，pdf\LaTeX{}
    无法编译）。连续编译两遍以生成目录与交叉引用：
    \begin{center}
      \texttt{xelatex main.tex}\quad\texttt{xelatex main.tex}
    \end{center}
    若本机使用中易字体，可将 \texttt{fontset=fandol} 改为 \texttt{fontset=windows}。
\end{enumerate}

% ---------------------------------------------------------------------------
% 目录
% ---------------------------------------------------------------------------
\clearpage
\tableofcontents
\clearpage

% ---------------------------------------------------------------------------
% 引言（原书 pp.1–2）
% ---------------------------------------------------------------------------
\section*{引言}
\addcontentsline{toc}{section}{引言}

同调镜像对称的核心，是以范畴论结构表达辛几何与代数几何之间的一种关系。仅这一句话，
就已道出本主题内在之美的许多，以及其雄心勃勃的目标所伴随的困难。

引导整个研究的这一猜想，由俄裔数学家 Maxim Kontsevich 在 1994 年于苏黎世举行的国际
数学家大会上的著名演讲（记录见 \cite{Kon94}）中首次提出；它以镜像对称的数值预言为
基础，而后者是从弦论中发展出来的一个数学分支。镜像对称断言一对"镜像"
Calabi--Yau 流形之间存在对偶，这类流形是某种特定的复辛流形。用物理的语言说，这些几何
对象具有弦振动据以发生的那些额外维数，而它们的对偶则表现为其中所描述的物理规律的
一致性。

Kontsevich 赋予镜像对称猜想以范畴的特征，用导出范畴的语言重新表述它——这是同调代数
极为关注的主题。粗略地说，所断言的对应该今只选取底层 Calabi--Yau 结构的一部分：即
给定一对镜像 Calabi--Yau 流形 $(X,X^\dagger)$，$X$ 的辛结构（"A-侧"，编码于其 Fukaya
范畴 $\mathcal F(X)$ 中）与 $X^\dagger$ 的代数几何结构（"B-侧"，编码于其凝聚层范畴
$\mathsf{Coh}(X^\dagger)$ 中）之间存在对偶；反之，$X$ 的 B-侧反映 $X^\dagger$ 的 A-侧。
具体而言，猜想存在两个三角范畴的等价：
\begin{equation*}
  \mathsf D^\pi(\mathcal F(X))\simeq \mathsf D^{\mathrm b}(\mathsf{Coh}(X^\dagger))
  \qquad\text{与}\qquad
  \mathsf D^\pi(\mathcal F(X^\dagger))\simeq \mathsf D^{\mathrm b}(\mathsf{Coh}(X)) .
\end{equation*}

我们的目标是尽可能合理地解释这些表达式的确切含义。事实上，这一图景如此丰富，我们只能
讨论 A-侧，对上述"更宏大的图景"仅作若干提示。此外，本文所提供的只是一个更为深刻、
因而也不可避免地更为艰深的理论的抛光版本：我们常常在显著简化的设定中工作，尤其是在
第二部分处理 Floer 理论时；只要合理，我们鼓励的是全面的概览，而非单一的深入剖析。
但愿这样的进路能传达同调镜像对称所借用的每一个纯数学分支的本质性。毕竟，不同数学领域
之间的这种协同，或许正是它最吸引人的优点。

全书分为两部分：对 $A_\infty$-范畴这一抽象代数概念作一般而自足的处理，以及随后对它
的一个具体实例——Fukaya 范畴——的专门化。具体地说，我们从第 1 章的
$A_\infty$-代数开始旅程，它们是 $A_\infty$-范畴的原型；第 2 章正式引入
$A_\infty$-范畴及必要的工具箱。第 3 章运用这套新技术锻造出与 $A_\infty$-范畴
$\mathcal A$ 相伴的有界导出范畴 $\mathsf D^{\mathrm b}(\mathcal A)$ 这一中间概念。
第 4 章把它丰富为同调镜像猜想中出现的分裂封闭导出范畴 $\mathsf D^\pi(\mathcal A)$。
接着，第 5 章把焦点转向辛流形，研究拉格朗日 Floer 同调——它对 Fukaya 范畴
$\mathcal F(X)$ 的定义不可或缺；第 6 章给出 $\mathcal F(X)$，并描述其
$A_\infty$ 结构与特殊的几何特征。最后，第 7 章综述导向同调镜像对称猜想的理论，
简要触及它的物理动机与在镜像对称中的起源，并力求对 B-侧也获得定性的理解。

\vspace{0.6em}
\noindent\textbf{致谢。}我要感谢我的导师 Will Merry，他仔细阅读了本论文的各个版本，
更重要的是，他始终支持我关于如何最好地发展论文的想法，直到它得以成文。
